\newpage
\section*{Questão 1 — Motor Otto 4T}
\label{sec:q_1}
\vspace{{10pt}}
\begin{{minipage}}{{\textwidth}}
\small
\paragraph{{Enunciado:}}
\begin{{verbatim}}
Um motor de ignição por compressão de 4 tempos apresenta potência de 
\textbf{62 kW} a 2200 rpm. A eficiência térmica ao freio é de 30\% e 
o poder calorífico inferior do combustível é de \textbf{42 MJ/kg}. 
O diâmetro do cilindro é \textbf{122.0 mm} e o curso 100 mm. 
Considerando densidade do ar = \textbf{1.184 kg/m$^3$}:

\begin{enumerate}
\item Determinar o consumo de combustível em kg/h;
\item Consumo de ar em m$^3$/h;  
\item Eficiência térmica indicada;
\item Eficiência volumétrica;
\item Pressão média efetiva;
\item Velocidade média do pistão.
\end{enumerate}

\end{{verbatim}}
\end{{minipage}}
\vspace{{15pt}}
\subsection*{Solução}
\noindent\textbf{Dados:}\\[5pt]
\begin{center}
\begin{tabular}{|c|r|c|} \hline
\textbf{Parâmetro} & \textbf{Valor} & \textbf{{Unidade}} \\ \hline\hline
$Pb$ & 62000.0 & W \\ \hline
$N$ & 2200.0 & min^{-1} \\ \hline
$eta_th_f$ & 0.3 & -- \\ \hline
$eta_m$ & 0.82 & -- \\ \hline
$PCI$ & 42000000.0 & J/kg \\ \hline
$AF$ & 15 & -- \\ \hline
$rho_ar$ & 1.184 & kg/m$^3$ \\ \hline
$d$ & 0.122 & m \\ \hline
$c$ & 0.1 & m \\ \hline
\end{tabular}
\end{center}\\[10pt]
\paragraph{Resolução:}
\begin{enumerate}[leftmargin=*,itemsep=8pt]
\subsubsection*{Consumo de combustível (kg/h)}
\noindent Fórmula padrão para motores: $\dot{{m}}_f = \frac{{P_e}}{\eta_{{t,e}} \cdot PCI} \times 3600$. Convertendo $P_e$ para W e isolando obtemos o consumo em kg/h.
\[\dot{m}_{f} = \frac{3600 Pb}{PCI \eta_{th f}}\]
\[\Rightarrow 17.714285714285715\]
\vspace{{6pt}}
\subsubsection*{Vazão volumétrica de ar (m$^3$/h)}
\noindent Aplicamos a fórmula $\dot{{V}}_{{ar}} = \frac{{AF_m \cdot \dot{{m}}_f}}{\rho_{{ar}}}$. Multiplicamos o consumo de combustível por hora pela razão ar/combustível e dividimos pela densidade do ar para obter a vazão em m$^3$/h.
\[V_{ar} = \frac{AF \dot{m}_{f}}{\rho_{ar}}\]
\[\Rightarrow 224.42084942084944\]
\vspace{{6pt}}
\subsubsection*{Eficiência térmica indicada}
\noindent Aplicamos a definição $\eta_{t,i} = \frac{\eta_{t,e}}{\eta_m}$, que relaciona a eficiência térmica ao freio com a eficiência indicada, considerando perdas mecânicas no motor.
\[\eta_{th i} = \frac{\eta_{th f}}{\eta_{m}}\]
\[\Rightarrow 0.36585365853658536\]
\vspace{{6pt}}
\subsubsection*{Cilindrada unitária (cc) (1 cilindro)}
\noindent Calculamos o volume deslocado por 1 cilindro usando a fórmula $V_d = \frac{\pi \cdot d^2}{4} \cdot c$, onde $d$ é o diâmetro interno e $c$ o curso do pistão.
\[V_{d} = \frac{\pi c d^{2}}{4}\]
\[\Rightarrow 0.001168986626400762\]
\vspace{{6pt}}
\subsubsection*{Eficiência volumétrica}
\noindent Para motor 4 tempos, a frequência de admissão por cilindro é $f_{{ciclo}} = \frac{{N}}{{120}}$. Calculamos a vazão volumétrica teórica de admissão por segundo como $\dot{{V}}_{{adm,teor}} = V_d \cdot f_{{ciclo}}$, considerando os 4 cilindros. Assim, a eficiência volumétrica é definida como $\eta_v = \frac{{\dot{{V}}_{{ar}}}}{{\dot{{V}}_{{adm,teor}}}}$.
\[0.0623391248391248 = 18.3333333333333 \pi c d^{2} \eta_{v}\]
\[\Rightarrow 0.7271930712298362\]
\vspace{{6pt}}
\subsubsection*{Pressão média efetiva (Pa)}
\noindent A pressão média efetiva é dada por $p_{{me}} = \frac{{P_e}}{{V_d \cdot f_{{ciclo}}}}$, onde $V_d$ é o volume total deslocado e $f_{{ciclo}} = \frac{{N}}{{120}}$ em motores 4 tempos.
\[pme = \frac{0.0545454545454545 Pb}{\pi c d^{2}}\]
\[\Rightarrow 723237.1409223459\]
\vspace{{6pt}}
\subsubsection*{Velocidade média do pistão (m/s)}
\noindent A velocidade média do pistão é calculada por $U_p = \frac{2 \cdot c \cdot N}{60}$, considerando dois deslocamentos lineares por rotação do virabrequim.
\[Up = \frac{N c}{30}\]
\[\Rightarrow 7.333333333333333\]
\vspace{{6pt}}
\end{{enumerate}}
\newpage